\documentclass[11pt]{article}
\usepackage[margin=1in]{geometry}
\usepackage{amsmath,amssymb}
\usepackage{graphicx}
\usepackage{booktabs}
\usepackage{hyperref}
\usepackage{natbib}

\title{Resolving the Mode of Action of Thiolutin: Complex In Vivo Effects but Direct RNA Polymerase II Inhibition In Vitro}
\author{Anonymous Authors}
\date{}

\begin{document}
\maketitle

\begin{abstract}
Thiolutin, a dithiolopyrrolone natural product widely used as a transcription inhibitor for mRNA stability studies, has an unresolved mechanism of action. Recent work attributed its activity entirely to zinc chelation, contradicting classic studies showing direct polymerase inhibition. Here we resolve this conflict through complementary genetic and biochemical approaches. Three independent genetic screens in \textit{Saccharomyces cerevisiae}---UV mutagenesis, Variomics library screening, and high-throughput Bar-seq of deletion and Variomics libraries---identify multidrug resistance, oxidative stress, metal homeostasis, and proteasome pathways as determinants of thiolutin sensitivity. Genetic dissection of the thioredoxin pathway reveals that thiolutin oxidizes thioredoxins (Trx1/Trx2) in vivo: the \textit{trx1$\Delta$ trx2$\Delta$} double deletion confers resistance and suppresses \textit{trr1$\Delta$} hypersensitivity. Thiolutin induces oxidative stress through redox cycling, depleting glutathione and driving nuclear localization of the redox-sensitive transcription factor Yap1. We discover that $\text{Mn}^{2+}$ is a required cofactor: thiolutin plus DTT plus $\text{Mn}^{2+}$ together---but not any pair alone---directly inhibit purified 12-subunit Pol~II transcription initiation in vitro. The active species must bind Pol~II before template DNA, exhibiting clamp-inhibitor-like behavior. Excess DTT abrogates inhibition, indicating that a specific partially reduced form is the active inhibitor. On transcription bubble templates that bypass the initiation block, pause-prone defective elongation is observed. Genome-wide Pol~II ChIP-seq after thiolutin treatment shows widespread occupancy changes consistent with Tor pathway inhibition rather than specific transcription targeting, indicating that in vivo transcriptional effects are predominantly indirect.
\end{abstract}

\section{Introduction}

Thiolutin is a dithiolopyrrolone natural product that has been widely used as a transcription inhibitor in mRNA stability studies for decades \citep{tipper1973action, tokunaga1973thiolutin}. Despite this extensive use, its mechanism of action remains contested. Classic studies demonstrated that thiolutin could inhibit partially purified RNA polymerase preparations, suggesting direct enzyme targeting. However, recent work showed that reduced thiolutin and its close relative holomycin chelate $\text{Zn}^{2+}$ and concluded that RNAP inhibition is indirect, mediated through metalloprotein inhibition \citep{chan2017reduced}.

This discrepancy has important implications: if thiolutin's transcriptional effects are indirect, then decades of mRNA stability measurements using thiolutin as a transcription shut-off tool may need reinterpretation. Several key gaps remain. No systematic genetic screen has identified the cellular pathways controlling thiolutin sensitivity. The roles of $\text{Mn}^{2+}$ and $\text{Cu}^{2+}$ in thiolutin activity have been largely overlooked. The specific reductants needed for thiolutin activation in vivo are unclear. And genome-wide Pol~II occupancy changes after thiolutin treatment have not been profiled.

\section{Results}

\subsection{Thiolutin Alone Does Not Inhibit Purified Pol~II}

Using purified 12-subunit yeast Pol~II with ssDNA templates, thiolutin at 10--25 $\mu$g/mL---with or without equimolar DTT---showed no detectable transcription inhibition across at least three independent experiments. This confirmed the recent negative result and motivated the search for missing cofactors.

\subsection{Three Genetic Screens Identify Thiolutin Sensitivity Pathways}

A UV mutagenesis forward genetic screen, manual Variomics library screening \citep{slusarczyk2023variomics}, and high-throughput Bar-seq of Variomics and deletion libraries identified convergent pathways. Gain-of-function mutations in multidrug resistance (MDR) transcription factors YAP1 and YRR1 and the efflux pump SNQ2 conferred thiolutin resistance. Loss of PDR1 or SNQ2 caused hypersensitivity, consistent with thiolutin being an efflux substrate. Thioredoxin reductase (TRR1) loss caused hypersensitivity, while TRR1 overexpression conferred resistance. Metal homeostasis genes, including ZAP1 ($\text{Zn}^{2+}$ homeostasis), and proteasome components were also identified.

\subsection{Thioredoxin Pathway Dissection}

Epistasis analysis revealed that \textit{trr1$\Delta$} hypersensitivity is due to accumulation of oxidized thioredoxins: the \textit{trx1$\Delta$ trx2$\Delta$} double deletion conferred resistance and suppressed \textit{trr1$\Delta$} hypersensitivity. This genetic evidence supports a model where thiolutin oxidizes thioredoxins in vivo.

\subsection{Thiolutin Induces Oxidative Stress via Redox Cycling}

Yap1-EGFP showed strong nuclear accumulation upon thiolutin treatment (10 $\mu$g/mL), comparable to $\text{H}_2\text{O}_2$ exposure \citep{herrero2008redox}. Glutathione measurements revealed significant depletion of total glutathione ($p \leq 0.01$, two-tailed paired $t$-test). DTT-reduced thiolutin could be reoxidized, consistent with redox cycling capacity.

\subsection{$\text{Mn}^{2+}$ Is the Missing Cofactor for Direct Pol~II Inhibition}

The critical discovery was that thiolutin + DTT + $\text{MnCl}_2$ together---but not any pair alone---directly inhibited Pol~II transcription initiation in vitro. This three-component requirement resolves the conflict between classic studies (which used partially purified fractions likely containing $\text{Mn}^{2+}$) and recent negative results (which used purified Pol~II without $\text{Mn}^{2+}$).

\subsection{Clamp-Inhibitor-Like Behavior}

Order-of-addition experiments showed that the thiolutin/$\text{Mn}^{2+}$/DTT complex must contact Pol~II before template DNA binding for effective inhibition. Pre-formed Pol~II-DNA complexes were resistant. This behavior is characteristic of RNAP clamp inhibitors \citep{belogurov2009structural}. Excess DTT abrogated inhibition, indicating that a specific partially reduced form---not fully reduced thiolutin---is the active species.

\subsection{Pause-Prone Elongation on Bubble Templates}

When the initiation block was bypassed using pre-melted bubble templates, thiolutin/$\text{Mn}^{2+}$/DTT produced defective, pause-prone elongation rather than complete block. This distinguishes thiolutin from classical clamp inhibitors that do not affect elongation on bubble templates.

\subsection{In Vivo Transcriptional Effects Are Predominantly Indirect}

Genome-wide Pol~II ChIP-seq after thiolutin treatment showed widespread occupancy changes. The dominant patterns---decreased occupancy at ribosomal protein genes and increased occupancy at stress-responsive genes---were consistent with Tor pathway inhibition \citep{loewith2002tor} and the environmental stress response, rather than direct, specific transcription targeting.

\section{Discussion}

Our results provide a unified model for thiolutin's mode of action. Upon cell entry (subject to MDR efflux), thiolutin's intramolecular disulfide is reduced by thioredoxins and/or glutathione. The reduced form chelates $\text{Zn}^{2+}$ (explaining metalloprotein effects) and undergoes redox cycling (generating oxidative stress). Critically, a thiolutin/$\text{Mn}^{2+}$/reductant complex directly inhibits Pol~II initiation. However, in vivo, the dominant transcriptional effects arise indirectly through Tor pathway inhibition and stress responses.

The practical implication is clear: thiolutin should not be used as a general transcription shut-off tool for mRNA stability studies, as its in vivo effects are dominated by indirect signaling perturbations rather than direct polymerase inhibition.

\section{Methods}

Forward genetic screens used UV mutagenesis of CKY457 followed by selection on 10 $\mu$g/mL thiolutin plates with bulk segregant analysis and whole-genome sequencing. Variomics and deletion library screens used Bar-seq with three biological replicates. In vitro transcription used purified 12-subunit yeast Pol~II with radiolabeled NTPs on ssDNA, dsDNA, and bubble templates. In vivo validation included Yap1-EGFP relocalization imaging, glutathione assays, serial dilution plate assays, growth curves, and Pol~II ChIP-seq.

\bibliographystyle{plainnat}
\bibliography{references}

\end{document}
